% Farben für jede Schicht (\providecommand: auch im Skript definiert)
\providecommand{\applicationLayerColor}{Gold}
\providecommand{\transportLayerColor}{Indigo}
\providecommand{\networkLayerColor}{RoyalBlue}
\providecommand{\dataLinkLayerColor}{Gray}

% command to make sure colorboxes don't create whitespace between code lines
\fboxsep=.6pt

% default settings for listings
\lstset{style=python}

\title[Netzwerke]{Netzwerke}

\subtitle{Transportschicht (Skript, Abschnitt 1.5)}

\author[Blum, Cyril] % (optional)
{Cyril Blum}

\institute[KLW] % (optional)
{
    Fachschaft Informatik\\
    Kantonsschule im Lee
}

\date[\today] % (optional)

\logo{\includegraphics[width=2cm]{Logos/FreeFlower.pdf}}

\titlegraphic{\vspace{-.7cm}\includegraphics[width=.85\textwidth]{"Grundlagen_Info/06_Datenbanken/Figures/title-emojis"}}

\frame{\titlepage}

\section{Transportschicht}
\setbeamercolor{frametitle}{bg=\transportLayerColor!20}

\begin{frame}{Netzwerk-Spiel: \textbf{Transportschicht}}
    Endgeräte trennen die Nachrichten in Segmente auf und fügen Sequenznummern hinzu:
    \begin{table}[H]
        \centering

        \begin{tabular}{p{2cm}l}
            \rowcolor{\transportLayerColor!20}Sequenz:                                    & 1                   \\\midrule
            \rowcolor{\applicationLayerColor!20}Nachricht\newline\tiny (max. 10 Zeichen): & \texttt{Hi Bob, wi} \\
        \end{tabular}
        \begin{tabular}{p{2cm}l}
            \rowcolor{\transportLayerColor!20}Sequenz:                                    & 2                  \\\midrule
            \rowcolor{\applicationLayerColor!20}Nachricht\newline\tiny (max. 10 Zeichen): & \texttt{e geht's?} \\
        \end{tabular}
    \end{table}
\end{frame}

\begin{frame}[fragile]{\faBoxOpen{} \acrshort{TCP} (\textit{Transmission Control Protocol}) Header}{Transportschicht}
    \acrshort{TCP} = Übertragungssteuerungsprotokoll
    \begin{table}[H]
        \centering
        \begin{tabular}{>{\raggedright\arraybackslash}p{4cm}>{\raggedright\arraybackslash}p{4cm}}

		\rowcolor{\transportLayerColor!30}\multicolumn{2}{c}{\textbf{\gls{TCP} Header}}       \\ \midrule
		\rowcolor{\transportLayerColor!15}Source Port & Destination Port                      \\ \midrule
		\rowcolor{\transportLayerColor!15}\multicolumn{2}{c}{Sequence \# (SEQ)}               \\ \midrule
		\rowcolor{\transportLayerColor!15}\multicolumn{2}{c}{Acknowledgement \# (ACK)}        \\ \midrule
		\rowcolor{\transportLayerColor!15}\multicolumn{2}{c}{Weitere \gls{TCP}-Header-Felder} \\ \midrule
		\rowcolor{orange!30}\multicolumn{2}{c}{\textbf{Nutzdaten}}                            \\ \midrule
		\rowcolor{orange!15}\multicolumn{2}{c}{$\ldots$ Anwendungsdaten $\ldots$}             \\
	\end{tabular}
    \end{table}
\end{frame}

\begin{frame}[fragile]{\faShip{} Ports}{Transportschicht}
    \begin{table}[H]
        \centering
        \footnotesize
        \begin{tabular}{ccc}
            \textbf{Portnummer} & \textbf{Anwendung} & \textbf{Verwendung}             \\
            \midrule
            20                  & \acrshort{FTP}                & Dateitransfer                   \\
            21                  & \acrshort{FTP}                & Befehlssteuerung                \\
            22                  & \acrshort{SSH}                & Sichere Shell-Zugriffe          \\
            25                  & \acrshort{SMTP}               & E-Mail-Versand                  \\
            53                  & \acrshort{DNS}                & Namensauflösung                 \\
            80                  & \acrshort{HTTP}               & Webseiten-Abruf                 \\
            110                 & \acrshort{POP3}               & E-Mail-Abholung                 \\
            143                 & \acrshort{IMAP}               & E-Mail-Management               \\
            443                 & \acrshort{HTTPS}              & Verschlüsselter Webseiten-Abruf \\
            993                 & \acrshort{IMAP}               & E-Mail-Management über \acrshort{TLS}      \\
            995                 & \acrshort{POP3}               & E-Mail-Abholung über \acrshort{TLS}        \\
        \end{tabular}
        \caption*{Gängige Ports, Anwendungen und deren Verwendung (\acrshort{DNS} meist über \acrshort{UDP}, sonst \acrshort{TCP})}
    \end{table}
\end{frame}

\begin{frame}{\faShip{} Ports}{Auftrag (Skript)}
    \begin{itemize}
        \item \aufgabebox{1.28 (Ports für \acrshort{HTTP}-Verbindungen)}
        \item \challengebox{1.29 (Weitere Port-Aufgaben)}
    \end{itemize}
\end{frame}

\begin{frame}[fragile]{\faHandshake{} \acrshort{TCP}: 3-way handshake}{Transportschicht}
    \begin{columns}[c]
        \begin{column}{.6\textwidth}
            \begin{figure}[H]
                \centering
                \adjustbox{max width=\textwidth, max height=.75\textheight}{
                    \begin{tikzpicture}[>=Stealth, node distance=2cm and 4cm]
		% Draw the entities
		\tikzset{c/.append style={fill=\transportLayerColor!30,rounded corners, align=center, text width=5cm}}
		\tikzset{person/.append style={text width=2.5cm,fill=ForestGreen!30,rounded corners,align=center}}
		\node[person](Alice) {\faFemale{} Alice\\\texttt{10.0.0.1}};
		\node[person, right=of Alice] (Bob) {\faMale{} Bob\\\texttt{10.0.0.2}};

        % Draw the arrows and labels
		\draw[->, thick, \transportLayerColor, dashed] (Alice.south) -- node[c] { Ich möchte mit dir sprechen, Bob. \\\texttt{SYN}, \texttt{SEQ} = $x$} ([yshift=-2.5cm]Bob.south);

		\draw[->, thick, \transportLayerColor, dashed] ([yshift=-2.5cm]Bob.south) -- node[c] { Ok, lass uns reden, Alice!\\\texttt{SYN, SEQ} = $y$ \texttt{ACK}=$x+1$} ([yshift=-4cm]Alice.south);
		\draw[->, thick, \transportLayerColor, dashed] ([yshift=-4cm]Alice.south) -- node[c] { Ok, danke Bob.\\\texttt{ACK}=$y+1$} ([yshift=-6cm]Bob.south);

		\node[below of=Alice, yshift=-5cm] (AliceZeit) {Zeit};
		\node[below of=Bob, yshift=-5cm] (BobZeit) {Zeit};

		% Time line
		\draw[->, thick] (Alice.south) -- node[xshift=-.2cm, rotate=270] {Zeitachse von Alice} (AliceZeit);
		\draw[->, thick] (Bob.south) -- node[xshift=.2cm, rotate=270] {Zeitachse von Bob} (BobZeit);
	\end{tikzpicture}
                    }
            \end{figure}
        \end{column}\vrule{}
        \begin{column}{.4\textwidth}
            \footnotesize
            \begin{enumerate}
                \item Computer A (\textit{client}, \texttt{10.0.0.1}) erstellt eine Verbindungsanfrage zum Computer B (\textit{server}, \texttt{10.0.0.2}) mittels eines Pakets, bei dem nur das \texttt{SYN}-Flag gesetzt ist.
                \item Der Server antwortet mit einem Paket, bei dem sowohl das \texttt{SYN}- als auch das \texttt{ACK}-Flag gesetzt sind
                \item Im letzten Schritt antwortet der \textit{client} nochmals mit einem einzelnen \texttt{ACK}-Flag\\
                      $\rightarrow$ \acrshort{TCP}-Verbindung ist nun erstellt!
            \end{enumerate}
        \end{column}
    \end{columns}
\end{frame}

\begin{frame}[fragile]{\faHandshake{} \acrshort{TCP}: 3-way handshake}{Auftrag}
    \begin{itemize}
        \item Skript: \aufgabebox{1.30-1.33 (\acrshort{TCP}, Server / Clients)}
        \item \challengebox{1.34 (E-Mail-Server)}
    \end{itemize}
\end{frame}


